\documentclass[11pt]{article}
\usepackage[a4paper,margin=25mm]{geometry}
\usepackage{amsmath,amssymb,bm}
\usepackage{graphicx}
\usepackage{booktabs}
\usepackage{microtype}
\usepackage{xcolor}
\usepackage{xr-hyper}
\usepackage[colorlinks=true,allcolors=blue]{hyperref}
\externaldocument[M-]{main}
\graphicspath{{figures/}}
\renewcommand{\thesection}{S\arabic{section}}
\renewcommand{\theequation}{S\arabic{equation}}
\renewcommand{\thetable}{S\arabic{table}}
\renewcommand{\thefigure}{S\arabic{figure}}
\newcommand{\todo}[1]{\textcolor{red}{[#1]}}
\newcommand{\dd}{\mathrm{d}}
\newcommand{\x}{\bm{x}}

\title{Supplement: Species volume fractions in granular mixtures from rotating radiography}
\author{}
\date{}

\begin{document}
\maketitle

\section{Annulus weights}\label{sec:window}

The periodogram of a windowed patch has expectation
$\mathbb{E}\hat P(\bm q)=\sum_{\bm q'}|\hat W(\bm q-\bm q')|^2\,\mathbb{E}P(|\bm q'|)/\sum|\hat W|^2$,
the true spectrum convolved with the power spectrum of the window. For an
isotropic $\mathbb{E}P$ the mean over the bins of annulus $b$ is a weighted
sum of $\mathbb{E}P$ at radial nodes $q_r$,
\begin{equation}
\mathbb{E}\hat P_b=\sum_r\omega_{br}\,\mathbb{E}P(q_r).
\end{equation}
We compute $|\hat W|^2$ by zero-padding the window to four times the patch
width, keep it within six bins of the origin, and for every bin of the
annulus distribute the weight $|\hat W(\bm q-\bm q')|^2$ over the radial
nodes by linear interpolation in $|\bm q'|$. The nodes are spaced at a
quarter of the bin width. The weights depend only on the window and the
patch size and are computed once.

Without these weights, the model at the mean wavenumber of each annulus
underestimates the measured annulus values near the zeros of the form
factors by up to about a third in the default band, because the measured
values cannot reach the zeros. With the weights, the measured and modelled
spectra of a uniform packing agree to within 3\% in the band.

The synthetic radiographs average $2\times2$ point samples per pixel instead
of integrating over the pixel. The transfer function of that average,
$\sin(qp/2)/[2\sin(qp/4)]$ per axis, differs from the sinc of a square
aperture by up to 16\% in power at $0.6$ of the Nyquist wavenumber, and the
synthetic spectra are divided by it instead.

\section{Windows}\label{sec:windows}

Table~\ref{tab:windows} compares spectral windows on the default data,
for meshes of 2 and 4 elements across the specimen (about 630 and 80 large
grains per element). The last column is the mean reduced $\chi^2$ of the
annuli on the 2-element mesh after the variance components are fitted; a
value close to one means the periodogram scatter alone accounts for the
residuals. Patch means, and hence $\phi$, do not depend on the window.

\begin{table}[h]
\centering
\caption{Errors in $c$ and $\phi_{\rm large}$ for different spectral windows
(seed 0).}
\label{tab:windows}
\begin{tabular}{lccccc}
\toprule
 & \multicolumn{2}{c}{2 elements} & \multicolumn{2}{c}{4 elements} & \\
Window & $c$ & $\phi_{\rm large}$ & $c$ & $\phi_{\rm large}$ & $\chi^2$ \\
\midrule
Hann & 0.023 & 0.0074 & 0.051 & 0.0179 & 1.55 \\
Tukey, taper fraction 0.5 & 0.019 & 0.0063 & 0.046 & 0.0161 & 1.65 \\
Tukey, taper fraction 0.25 & 0.029 & 0.0084 & 0.055 & 0.0177 & 1.67 \\
Sine multitaper, $2\times2$ & 0.028 & 0.0092 & 0.052 & 0.0179 & 1.13 \\
Sine multitaper, $3\times3$ & 0.041 & 0.0126 & 0.064 & 0.0221 & 1.03 \\
Hann, stride 12 & 0.023 & 0.0064 & 0.052 & 0.0168 & 1.53 \\
\bottomrule
\end{tabular}
\end{table}

\section{Error model alternatives}

Two simpler weightings were tried before the one in the main text. Giving
each patch mean a fixed relative error weights the patches at the edge of
the specimen most heavily, where the chords are short and the relative
grain-count fluctuation is largest. This was not isolated in a controlled
comparison, because the same runs also contained the wall layer of
Section~\ref{M-sec:wall}. Leaving the detector noise out of the variance of
an annulus gives annuli whose signal is comparable to the noise a variance
near zero once the noise is subtracted. With the self-term kernel and a
noise of 3 voxels of path length per pixel, the error in $c$ was above 1
without the noise term and below 0.1 with it (errors at the mesh nodes).

\section{Reference fields}\label{sec:reference}

The realised field of species $a$ is the least-squares fit of
Eq.~\eqref{M-eq:fields} to the grains,
$\min\int(w_a\tilde\phi_a-I_a)^2\dd V$, with $I_a$ the indicator of
species $a$. Each grain is counted as its volume at its centre, so the
right-hand side is $\sum_j v_a w_a(\x_j)\bm N(\x_j)$ over the grains $j$ of
species $a$ in each packing, weighted by the time basis at the packing's
time. The mass matrix is integrated on a grid of spacing 3 voxels. A
gradient penalty of $10^{-4}$ of the mean diagonal keeps nodes whose support
holds little material finite.

The reference must be fitted on a basis shared by the reconstructions that
are compared with it. A reconstruction with three time knots, compared with
a reference fitted with three knots, appears accurate for a change over a
quarter turn, because neither can represent the change; against a reference
on 13 knots its error is 0.041 (Table~\ref{M-tab:time}). Likewise, a mesh
with one element across the specimen is bilinear across it and represents
any linear lateral gradient, so the layer-wise fields of
Table~\ref{M-tab:layers} tie all nodes of a layer together.

\section{Counting fluctuation of a projected field}\label{sec:floor}

Let $\tilde\phi_a(\x)=\bm N(\x)^\top\mathsf M^{-1}\bm b$ with
$\mathsf M=\int\bm N\bm N^\top\dd V$ over the specimen and
$\bm b=\sum_jv_a\bm N(\x_j)$ over the grains $j$ of species $a$. For grain
centres forming a stationary point process of density $\rho_a$ with
structure factor $S_{aa}$, and shape functions that vary slowly on the scale
of the grain spacing,
$\operatorname{cov}\bm b=\rho_av_a^2S_{aa}(0)\,\mathsf M$. Hence
\begin{equation}
\operatorname{var}\tilde\phi_a(\x)=\rho_av_a^2S_{aa}(0)\,\bm N^\top\mathsf M^{-1}\bm N
=S_{aa}(0)\,\phi_av_a\,\bm N^\top\mathsf M^{-1}\bm N,
\end{equation}
which is Eq.~\eqref{M-eq:floor}. For the solid fraction the variance is
$\sum_{ab}S_{ab}(0)\sqrt{\phi_av_a\phi_bv_b}\,\bm N^\top\mathsf M^{-1}\bm N$.

\section{Code checks}

The kernel of the patch means agrees with an independent ray-marching
projector of a smooth field to 2\% in the median over patches. The exact
projector of spheres reproduces the chord through the centre of a sphere of
diameter 10 voxels to 0.1 voxels and its volume to 1\%. In the dilute limit
the interference kernel of Eq.~\eqref{M-eq:chia} reduces to the self term
to 1\%, and the self term averaged over many lobes of the form factor
matches the Porod limit to 2\%. These checks are part of the test suite of
the code.

\end{document}
